\section{Analisi tecnica: funzionamento e debolezze sfruttate}
\label{sec:analisi-tecnica}

% Persone 2 e 3 — circa 1,6 pagine complessive, figure incluse.
\subsection{npm: dall'account compromesso alla dipendenza malevola}
% Persona 2 — circa 0,8 pagina.
% Spiegare il phishing, l'abuso dell'account publisher e il percorso dalla
% release malevola alla dipendenza e al bundle browser. Usare debug@4.4.2
% come esempio circoscritto; distinguere browser da server/CLI.

\subsection{AUR: dalla modifica della ricetta all'esecuzione}
\label{subsec:aur-tecnica}
% Persona 3 — circa 0,8 pagina.
\subsubsection{Abuso della manutenzione della ricetta}
Nel caso AUR, la fiducia non riguarda soltanto il programma finale: l'utente
scarica una ricetta di compilazione e si affida alle istruzioni che questa
contiene. L'AUR raccoglie ricette \texttt{PKGBUILD} inviate dagli utenti; non è
un repository binario ufficiale e le ricette non sono sottoposte a una verifica
completa da parte del progetto Arch~\cite{archwiki-aur}. Per ottenere un
pacchetto, \texttt{makepkg} interpreta il \texttt{PKGBUILD}, che può scaricare
file ed eseguire comandi durante la preparazione e la compilazione
~\cite{archwiki-makepkg}.

Un esempio circoscritto è \texttt{gnome-randr-rust}. L'11 giugno 2026, un
utente segnalò sulla lista ufficiale \texttt{aur-general} che un nuovo
manutentore aveva adottato il pacchetto e aggiunto uno script d'installazione
con il comando \texttt{npm install atomic-lockfile yargs}; il pacchetto non
aveva una ragione tecnica apparente per introdurre npm. Un manutentore AUR
confermò la natura malevola della modifica e ripristinò il commit precedente
~\cite{aur-mail-gnome-randr}. Il passaggio di fiducia abusato era la facoltà
di adottare e modificare una ricetta rimasta senza manutentore attivo, facendo
apparire legittime nuove istruzioni di installazione. Si tratta di un abuso
delle procedure di manutenzione AUR, non della compromissione del registro npm
né dei repository ufficiali di Arch.

La discussione tecnica sulla lista chiarì inoltre che il pacchetto npm
richiamato conteneva un hook di ciclo di vita \texttt{preinstall}, che avviava
un eseguibile Linux incluso nel pacchetto. In questo modo il comando aggiunto
alla ricetta non si limitava a scaricare una libreria: passava l'esecuzione al
codice distribuito attraverso il secondo ecosistema di pacchetti
~\cite{aur-mail-gnome-randr}. Il fatto che un pacchetto AUR introduca npm o un
altro gestore di pacchetti non prova da solo un attacco, ma quando la nuova
dipendenza non ha relazione con il software e avvia script durante
l'installazione costituisce un'anomalia da esaminare.

\subsubsection{Ambiente di esecuzione e privilegi}
La catena è: adozione del pacchetto $\rightarrow$ modifica della ricetta/script
\texttt{.install} $\rightarrow$ esecuzione di \texttt{npm install} nel callback
\texttt{post\_install}. La compilazione con \texttt{makepkg} e l'installazione
con \texttt{pacman} sono fasi distinte: Arch vieta di eseguire \texttt{makepkg}
come amministratore, mentre il callback di uno script \texttt{.install} viene
eseguito durante la transazione di installazione o aggiornamento
~\cite{archwiki-makepkg,arch-alpm-scriptlet}. Nel caso segnalato, quindi, il
codice malevolo si attivava se l'utente installava o aggiornava quel pacchetto;
il callback viene eseguito nel contesto della transazione di pacman e può
quindi ereditare i privilegi elevati usati per installare il pacchetto. Questo
non significa che il codice ottenga sempre privilegi di amministratore in
qualsiasi configurazione: conta la fase che lo avvia e l'eventuale separazione
dei privilegi adottata dall'utente. La fonte documenta il meccanismo, ma non
dimostra quante persone abbiano eseguito il pacchetto o quali sistemi siano
stati compromessi.

L'avviso Arch del 12 giugno conferma un'ondata di adozioni e aggiornamenti
malevoli e invita a controllare le modifiche a \texttt{PKGBUILD} e agli script
d'installazione~\cite{arch-aur-incident-2026}. Nel caso specifico, i
manutentori ripristinarono la revisione precedente e sospesero l'account
responsabile~\cite{aur-mail-gnome-randr}. I download o i pacchetti coinvolti
non equivalgono a installazioni riuscite né a vittime accertate.
